% Hand-written front matter (not generated). Edit freely.
\begin{titlepage}
  \centering
  \vspace*{2.5cm}
  {\Huge\bfseries\color{brand} \ProjectName\par}
  \vspace{0.6cm}
  {\LARGE \DocTitle\par}
  \vspace{0.4cm}
  {\large\itshape \DocSubtitle\par}
  \vspace{2.2cm}
  \begin{tabular}{@{}r@{\hspace{1em}}l@{}}
    \textbf{Team:} & Amine Essid \\
    \textbf{Product Owner:} & Amine Essid \\
    \textbf{Scrum Master:} & Amine Essid \\
    \textbf{Developers:} & Amine Essid \\
    \textbf{Organization:} & Amine Essid \\
    \textbf{Version:} & 1.0.1 \\
    \textbf{Date:} & \DocDate \\
  \end{tabular}
  \vfill
  {\small Scrum Guide terminology is used throughout: Product Backlog, Sprint Backlog, Increment, Product Goal and Sprint Goal.\par}
\end{titlepage}

\chapter*{About this document}
\addcontentsline{toc}{chapter}{About this document}
\markboth{About this document}{About this document}

This document collects the Scrum artifacts for the \textbf{\ProjectName}: an AI-native customer support platform that handles full conversations over text and voice, answers from approved organizational knowledge, performs explicitly authorized business actions, and hands over to human agents with context.

\paragraph{How it is organized.}
\begin{itemize}
  \item \textbf{Part I -- Product and Scrum Framework:} vision and scope, the Scrum charter and events, and an overview of the Product Backlog.
  \item \textbf{Part II -- Epics and User Stories:} the 12 epics and 38 user stories, each with its acceptance criteria and the diagrams that explain the behavior.
  \item \textbf{Part III -- Release and Sprint Planning:} release goals and risks, the provisional Sprint plan with one diagram per Sprint, and the Definitions of Ready and Done.
  \item \textbf{Part IV -- Architecture and Design Reference:} system architecture, module catalog, communication contracts, domain model, AI engine, voice, integrations, security and observability.
  \item \textbf{Part V -- Technology Report:} the proposed technology stack, the technologies for every epic, user story and Sprint, the role and usage rule of each technology, and the decisions still to validate.
  \item \textbf{Part VI -- UML Models:} use case, class, sequence, activity and structural (component, package, deployment, object) diagrams.
\end{itemize}

\paragraph{Status of the content.}
The backlog, sprint sequence and estimates are a \emph{planning hypothesis}, not a commitment. Team size, capacity and calendar dates are unknown; story points are team-specific and are left as ``TBD'' in the backlog table. Dates in the Gantt chart are placeholders. Names on the title page are placeholders to be replaced.

\paragraph{Conventions.}
\code{MUST} means required for correctness or security; \code{SHOULD} means recommended unless a documented reason exists. AI output is treated as untrusted input: authorization and business invariants are enforced by application code.

\paragraph{Diagram sources.}
All diagrams are Mermaid sources kept in \code{diagrams/mermaid/mmd/} and rendered to PDF by the Mermaid CLI. The Markdown files in \code{diagrams/mermaid/} are the editable originals.
